% ============================================================================
% October 9, 2026 build (agent draft). Policy section, new. Sources:
%   elastic stock (the deck's "Policy Results"):
%     output/model/policy_transition_battery_20261008/h100_torch/
%     (table_h100.md, tax_vs_credit_h100.pdf), base 14.402, from the 2023 state;
%   slow stock (deck): tax_vs_credit_h100_slow_stock.pdf;
%   by supply regime (deck appendix): ptax_by_supply_regime.pdf, h16 runs
%     (table_h16-20261008T153908-0eb8.md);
%   slow-stock refit and today's arms: output/model/slow_stock_h100_20261009/
%     README.md (project copy README-20261009T203735-9f66.md), EXPERIMENTAL;
%   channel reading: notes/ptax_vs_credit_channels_20261009.md;
%   long run: \pol macros from latex/JMP_slides/prod_tables/policy.tex
%     (reconciliation_14p40_20261007, stationary GE at psi 0.110).
% ============================================================================
\providecommand{\mocktodo}[1]{\textcolor{red}{[TODO: #1]}}

\section{A Housing Policy}
\label{sec:policy}

Section~\ref{sec:quant-2023} showed that large homes sit with older
households, and Section~\ref{sec:results} that the per-period cost of space is
what moves the number of children. A natural policy question follows: can a
policy that changes who holds housing, and at what price, raise births? I
study the experiment of \citet{coven2024}. They tax housing by value and
rebate the revenue equally to all households. Older owners, who hold the most
housing, pay more than they get back, and young renters get back more than
they pay. House prices fall, and constrained young households buy sooner or
buy bigger. I run the same experiment from the 2023 economy: the property tax
rises unexpectedly and permanently from 1.06 to 2 percent of value a year, and
its revenue is rebated equally to all households at every date. I follow
prices, rents, ownership and births along the path, and compare the tax with
the mortgage credit of Section~\ref{sec:results-2023}, raising the financed
share from 80 to 95 percent, also in general equilibrium. In both cases the
preference for children stays at its 2023 value, and every change is measured
against the baseline path from the same 2023 state.

\subsection{Results with the baseline housing supply}
\label{sec:policy-baseline}

Figure~\ref{fig:policy-baseline} reports the two policies when the housing
stock moves along the supply curve at every date. The tax lowers the house
price by 6.3 percent on impact, a gap that narrows slowly over the century.
Births fall by 0.7 percent in 2023 and are 0.3 to 0.5 percent above the
baseline from the 2030s through the 2060s. Ownership of the young rises by
about 3 percentage points on impact and stays about 2 points higher through
the 2060s. The adult population is 0.4 percent larger by 2100. Mortgage
credit raises the ownership of the young by 20 to 25 percentage points
throughout and barely moves the price. Births rise by 0.6 percent on impact,
are below the baseline from the late 2020s, and are about 2 percent lower
around 2100, when the population is 0.7 percent smaller.
% [MODEL] table_h100.md, elastic arm: price -6.33 (2023) to -6.23 (2063),
% -5.35 at the last date; births -0.69, +0.01, +0.25, +0.38, +0.47, +0.48,
% +0.68 at the last date; own 18-29 +2.83, 1.87, 1.73, 2.60, 2.54, 1.95;
% adults +0.67 at the long-run endpoint. LTV: own 18-29 +19.9 then 24.7-25.2;
% births +0.56, -0.26, -0.63, -0.64, -0.32, -0.81; adults -1.09 in the long
% run. The 2100 values are read from the figure.

\begin{figure}[t]
\centering
\includegraphics[width=0.95\linewidth,trim=0 13 0 0,clip]{../../output/model/policy_transition_battery_20261008/h100_torch/tax_vs_credit_h100.pdf}
\caption{Property tax and mortgage credit from 2023}
\label{fig:policy-baseline}
\par\smallskip
\begin{minipage}{0.92\linewidth}
\footnotesize\textit{Notes:} Change relative to the baseline path from 2023
in births, the house price, ownership at ages 18--29 and the adult population.
Property tax from 1.06 to 2 percent a year with an equal rebate; mortgage
credit: financed share from 80 to 95 percent. Housing stock on the supply
curve at every date.
\end{minipage}
\end{figure}

Why is the effect of the tax so small? A tax on value raises the cost of
holding a room. When the stock can shrink within a period, builders withdraw
rooms, so the price falls by less than the tax and the rent of a room rises:
by 13.5 percent on impact and by 10 to 12 percent later. The rebate is then
largely offset by a higher rent on the space a family needs, and births rise
by less than half a percent. At fixed prices, the same tax raises rent by 19.6
percent and lowers births by 2.9 percent on impact. What makes the tax
pro-natal, to the extent it is, is the fall in the price.

\subsection{A slowly adjusting housing stock}
\label{sec:policy-slow}

Housing is durable, and the stock cannot shrink by four percent in four
years. I therefore replace the baseline supply with a stock that adjusts
slowly. Builders replace depreciated rooms each period, and build more when
prices are high:
\begin{equation}
 I_t=\delta_H\,H_0\,P_t^{\eta},\qquad H_t=(1-\delta_H)H_{t-1}+I_t .
 \label{eq:policy-slow}
\end{equation}
In a steady state $I=\delta_H H$, so $H=H_0P^{\eta}$: the long-run supply
curve, the 2007 steady state and the calibration are unchanged, and only the
speed of adjustment differs. The stock moves toward the supply curve, closing
a fraction $\delta_H$ of the gap each year. I compare it with two extremes:
the baseline, in which the stock moves to the supply curve at once, and a
stock frozen at its 2023 level.

Figure~\ref{fig:policy-slow} repeats the two experiments with the slowly
adjusting stock. The tax now lowers the price by 11 percent on impact, and
the price gap closes over the century as the stock shrinks. Births rise by
about 2.5 percent on impact and by 3.5 to 4 percent through the middle of the
century, and the adult population is 3.5 percent larger by 2100. The
ownership of the young rises by 5 points on impact and is back to the
baseline by the 2040s. Mortgage credit is almost unaffected by the speed of
supply: young ownership rises by 19 to 27 points, births rise by less than
one percent at first and fall by 1.5 percent by 2100.

\begin{figure}[t]
\centering
\includegraphics[width=0.95\linewidth]{../../output/model/policy_transition_battery_20261008/h100_torch/tax_vs_credit_h100_slow_stock.pdf}
\caption{Property tax and mortgage credit with a slowly adjusting stock}
\label{fig:policy-slow}
\par\smallskip
\begin{minipage}{0.92\linewidth}
\footnotesize\textit{Notes:} As Figure~\ref{fig:policy-baseline}, with the
housing stock following \eqref{eq:policy-slow} from 2023.
\end{minipage}
\end{figure}

Figure~\ref{fig:policy-supply} isolates the role of supply for the tax. With
the baseline supply, rooms are withdrawn, the rent of a room rises by 13
percent, and births barely move. With a slow or a fixed stock, the rooms
stay, the house price falls instead, the rent barely moves on impact, and
births rise by about 3 percent on impact, and later by 3 to 4 percent with
the slow stock and by up to 6 percent with the fixed stock. The tax is
pro-natal when the price absorbs it, and nearly neutral when supply passes it
into rents.
% [MODEL] table_h16: frozen stock, births +2.85 on impact rising to +6.0,
% price -12.86, rent -0.10; elastic rent +13.55; fixed prices rent +19.64,
% births -2.93.
\mocktodo{the legends of the policy figures label the tax 1\%$\to$2\%; the
runs raise it from 1.06 to 2 percent. Relabel the figures.}

\begin{figure}[t]
\centering
\includegraphics[width=\linewidth]{../../output/model/jmp_slides_transition_figures_20261009/ptax_by_supply_regime.pdf}
\caption{The property tax under three housing supplies}
\label{fig:policy-supply}
\par\smallskip
\begin{minipage}{0.92\linewidth}
\footnotesize\textit{Notes:} Property tax from 1.06 to 2 percent a year,
rebated, from 2023: change relative to the no-tax path in births, the house
price, rent per room and the population, with the stock on the supply curve
at every date, adjusting slowly as in \eqref{eq:policy-slow}, or fixed at its
2023 level. Sixteen-period runs.
\end{minipage}
\end{figure}

\subsection{Where the effect of the tax comes from}
\label{sec:policy-channels}

The slowly adjusting stock changes the transition of
Section~\ref{sec:quant-transition}, so I re-estimate the two preference
shocks under it, matching the same two fertility rates, and run the tax from
the 2023 state of the re-estimated path.%
\footnote{The re-estimated shocks are $\psi_1=0.1295$ and $\psi_2=0.1062$,
which match period fertility of 1.861 and 1.646 to within 0.0004. The 2023
price is 10 percent below its 2007 level.}
Table~\ref{tab:policy-channels} reports the tax under four variants, against
the no-policy path from the same state. Births are reported in 2023, as the
mean over 2031--2035, and in 2067; completed fertility is that of the cohort
aged 18 to 21 in 2023.

\begin{table}[t]
\centering
\caption{Where the effect of the property tax comes from}
\label{tab:policy-channels}
\small
\begin{tabular}{lrrrrrr}
\toprule
 & \multicolumn{3}{c}{Births (\%)} & Completed & Price & Rent per \\
\cmidrule(lr){2-4}
 & 2023 & 2031--35 & 2067 & fertility (\%) & 2023 (\%) & room 2023 (\%) \\
\midrule
\emph{Slow stock, $\eta=0.63$} &&&&&& \\
\quad Tax, rebated & $+2.4$ & $+3.3$ & $+3.4$ & $+3.9$ & $-11.1$ & $+1.6$ \\
\quad Tax, rebate held at its base path & $-2.2$ & $-1.1$ & $-4.1$ & $-0.5$ & $-15.8$ & $-5.4$ \\
\quad Financed share $80\%\to95\%$ & $+0.9$ & $+0.2$ & $-1.1$ & $-0.6$ & $+0.1$ & $+1.8$ \\
\addlinespace
\emph{Slow stock, $\eta=1.5$} &&&&&& \\
\quad Tax, rebated & $+2.0$ & $+2.4$ & $+1.3$ & $+3.0$ & $-9.8$ & $+2.7$ \\
\addlinespace
\emph{Stock on the supply curve, $\eta=1.5$} &&&&&& \\
\quad Tax, rebated & $-2.0$ & $-1.7$ & $-3.1$ & $-1.8$ & $-3.6$ & $+18.0$ \\
\quad Financed share $80\%\to95\%$ & $+0.4$ & $-0.6$ & $-1.0$ & $-1.4$ & $+0.0$ & $+1.0$ \\
\bottomrule
\end{tabular}
\par\smallskip
\begin{minipage}{0.95\linewidth}
\footnotesize\textit{Notes:} Each row: change relative to the no-policy path
from the 2023 state of the transition re-estimated under the stated supply.
Tax: property tax from 1.06 to 2 percent a year. Completed fertility: sum of
the age-specific birth rates of the cohort aged 18--21 in 2023. With $\eta=1.5$
the supply scale is renormalized so that the 2007 steady state is unchanged.
\mocktodo{these runs are experimental (Oct 9): the slow-stock rule and
$\eta=1.5$ are not part of the calibrated model, the terminal check fails near
2400, and the $\eta=1.5$ stock-on-curve row starts from the closest grid pair
rather than the exact fit. The $\eta=1.5$ slow-stock credit arm is still
running.}
\end{minipage}
\end{table}

Three results stand out. First, the rebate carries the effect. When the
rebate is held at its base path, so that the extra revenue is not returned,
the same tax lowers births by 1.1 percent in 2031--2035 and completed
fertility by 0.5 percent, even though the price falls by 16 percent and the
rent of a room falls by 5 percent. The difference between the rebated and the
held tax is 4.4 points of births. Second, the sign of the effect follows how
much of the tax reaches the rent of a room. With a slow stock, the price
absorbs the tax and the rent rises by 1.6 to 2.7 percent on impact; births
rise by 2 to 3 percent. With a stock that can shrink quickly and an elastic
supply, the price falls by only 3.6 percent, the rent rises by 18 percent, and
births fall by about 2 percent, as at fixed prices. Third, mortgage credit
raises the ownership of the young by 19 to 21 percentage points on impact and
keeps births within about one percent of the baseline at every date, and
completed fertility within 1.5 percent, under both supplies for which the
credit arm has run.

I read these results as follows. When the price absorbs the tax, the rebated
tax is a transfer to every household, paid for by a one-time capital loss on
the households that own in 2023, at an almost unchanged rent per room. The
transfer raises the buffer of the households that decide on a first child
without savings, the margin on which Section~\ref{sec:results} found births
most responsive; the price fall itself lowers the one-time cost of buying,
which in Section~\ref{sec:results-2023} moved the timing of births and not
their number. When the price does not absorb the tax, it is a tax on rooms,
partly offset by a flat rebate, and a tax on rooms is a tax on the space
children need. The durability of the stock decides which case applies, and it
also decides who gains. The cohort young at the announcement receives the
rebate at an almost unchanged rent and has 3.9 percent more children; the
cohort aged 18 to 21 in 2031 gains less, 2.8 percent, consistent with a stock
that shrinks over time and a rent that rises with it. Births are counts, so
part of their rise at later dates reflects the larger cohorts born after the
reform rather than more children per household. The burden falls on the
owners of 2023. In the model they are less concentrated among the old than in
the data, because young families hold twice their share of the large homes
(Figure~\ref{fig:quant-allocation}).
% [REVIEW Oct 10] notes/ptax_vs_credit_channels_20261009.md marks the
% buffer and timing readings as inference [I]; the cohort numbers (3.9, 2.8)
% are in notes/housing_tenure_claims_20261009.md; "the owners are mostly old"
% is contradicted by notes/argument_for_corina_20261008.md step 17 (old
% childless hold 21.6% of large homes in the model against 40% in the data).

Two qualifications matter for how far this reading goes. The tax does not
move space to young families. Rent per room does not fall in any rebated arm,
and the gain in young ownership is small and fades. What reaches young
families is income, not rooms. And it remains open whether the housing
market matters for delivering that income. A transfer of the same size
financed by a payroll tax would show whether it matters that the rebate is
funded by a tax on housing, which falls largely on retired owners and on a
capitalized price rather than on the earnings of would-be parents.
\mocktodo{the payroll-financed arm is specified
(notes/ptax\_vs\_credit\_channels\_20261009.md) and waits on approval of the
earnings-tax patch.}

\subsection{The long run}
\label{sec:policy-longrun}

Because the population replaces itself only at 2.1 children per household,
each policy returns births per household to replacement in the long run; its
lasting effects appear in the population and in prices.
Figure~\ref{fig:policy-longrun} compares stationary equilibria at the 2023
preference for children, with and without each policy, under the baseline
supply. Relative to the stationary equilibrium without it, the property tax
changes the house price by $\pol{end_ptax2.price_pct.1}$ percent, the
population by $\pol{end_ptax2.pop_pct.1}$ percent and young ownership by
$\pol{end_ptax2.own_18_29_pp.1}$ points. Mortgage credit changes young
ownership by $\pol{end_phi095.own_18_29_pp.1}$ points, the price by
$\pol{end_phi095.price_pct.1}$ percent and the population by
$\pol{end_phi095.pop_pct.1}$ percent. At fixed prices, the tax would change
births by $\pol{ptax2_fp.births_pct.1}$ percent: the price fall is what makes
it pro-natal.

\begin{figure}[t]
\centering
\includegraphics[width=\linewidth]{../../output/model/jmp_slides_policy_functions_20261007/longrun_bars.pdf}
\caption{Property tax and mortgage credit in the long run}
\label{fig:policy-longrun}
\par\smallskip
\begin{minipage}{0.92\linewidth}
\footnotesize\textit{Notes:} Stationary equilibrium at $\psi_2=0.110$ with
each policy, relative to the stationary equilibrium without it: house price,
population and ownership at ages 18--29. Baseline housing supply.
\end{minipage}
\end{figure}
